\section{Conclusions and Insights}
\label{sec:conclusions}

% ~0.75 page. Managerial takeaways, tied to the KPIs.
\begin{itemize}\setlength\itemsep{1pt}
  \item The primary constraint is \textbf{drying} (highest required unit-hours
        and \SI{94.9}{\percent} simulated utilisation), with the
        \textbf{cleaning team} second (\SI{79.6}{\percent}); packaging, despite
        the largest absolute workload, is third (\SI{52.2}{\percent}) because it
        runs on three parallel lines.
  \item \textbf{No capacity lever raises throughput.} A second dryer relieves
        drying but shifts the bottleneck onto the cleaning team
        (utilisation \SI{79.6}{\percent}$\rightarrow$\SI{94.8}{\percent}), since
        each dryer adds a 2-hour clean; adding cleaning staff then shifts it
        onto packaging. Shipped output holds at 60{,}000--62{,}000 bottles
        throughout, because volume is fixed by the 50-batch plan and per-batch
        yield, not by station capacity. The $2^3$ factorial confirms this: of
        the seven effects, only a fourth packaging line is significant, and only
        marginally ($+717$ bottles, $+\SI{1.2}{\percent}$).
  \item Resequencing by product has no measurable effect on throughput
        ($-\SI{0.6}{\percent}$, not significant); its only benefit is a small
        reduction in flow time.
  \item Towards 2030, the line absorbs demand compressed up to $\times2.0$ with
        \textbf{zero release violations}; congestion rises (flow time
        $+\SI{13}{\percent}$, WIP $+\SI{14}{\percent}$) but no unbounded backlog
        forms. Drying --- not the schedule --- is the resource that would bind
        first if the plan's batch volume itself grew.
  \item \textbf{Industry~4.0 is the decisive lever.} Cutting press scrap from
        \SI{10}{\percent} to \SI{4}{\percent} raises shipped throughput by
        4{,}764 bottles per run ($+\SI{7.8}{\percent}$, 95\% CI
        $[3{,}479,\,6{,}050]$) --- over six times the only significant
        capacity effect, achieved through a one-off sensor retrofit with no added
        line capacity and no recurring cost.
\end{itemize}
